\section{The Purpose of Life}

\begin{quotex}
We experience the reality of God not as something indefinitely divine, but we are conscious of Him as He really is, all-perfect or absolute. And our soul too is revealed to us in our inner experience not merely as something distinct from material facts, but as a positive force which struggles with the material processes and overcomes them.
\flright{\textsc{Vladimir Solovyov}}
\end{quotex}


It would be quite strange in our time for people to be wondering how to make fire or use a wheel. Some things are simply settled. Yet, a great number of people are still searching for the ``purpose of life'', wandering aimlessly like a lion on the savannah. And most of those who are not searching, most likely have not even given it a thought.

This is despite millennia of religious, philosophical, and scientific movements, each of which makes a claim to reveal the true purpose of life. This smorgasbord of options is hardly a boon to the many seekers, but rather a burden. It would be impossible in a lifetime to study each particular religious tradition, along with its various aberrations and variations, in order to arrive at some ``objective'' conclusion. This is not merely an intellectual problem, but more importantly, a religious tradition can only truly be known from the ``inside'', i.e., by participation. Paradoxically, then, a seeker must make a commitment prior to any confirmation.

The end result is a battle of worldviews, each claiming exclusivity. When restricted to the intellectual plane, the infinite conversation or debate arises, with no possibility of resolution. Undeterred, those who choose to engage in that endless chitchat, hope beyond hope to find the ultimate ``defeater'', the \emph{bon mot}, that will force the rival to surrender. Hence, there is the unending drama of conversion, apostasy, heresy, and back again. One's identity becomes raveled in an external thought form, with its highest expression in code words and intellectual ``tribalism'', or even facebook ``likes''.

The esoterist, however, having passed through the first trial, has given up all such thought forms. With the consciousness emptied of all clutter, the esoterist sees that objective truth can be found in one's subjectivity. Hence, the task no longer takes the form of an intellectual debate, but rather that of true intellectuality based on intuition, and is more like ``seeing'' than ``thinking''. We will see then that:

\textit{The purpose of life is realized in the process of manifesting God in matter.}

\paragraph{The Reality of the Moral Order}


\textbf{Vladimir Solovyov}, in \emph{The Justification of the Good}, developed a useful schema for this task. Note, first of all, that this schema is independent of any particular thought form or worldview. Rather, it is something that can be observed in one's consciousness and each kingdom, or stage, represents a particular state of being. This is a very rich representation, which I am integrating with a Traditional metaphysic.

This involves the moral order, since it depends on action, and not on passive contemplation. As the motto above shows, Solovyov accepts \textbf{Giambattista Vico}'s \emph{verum factum} principle: truth is known through creation, not through observation. In short, it requires an act of will. After all, who knows the way to San Jose: one who has contemplated it on a map or the one who has actually made the trip?

Solovyov identifies the traditional five kingdoms, as shown in table \ref{tab:FiveStatesBeing}. Associated with each kingdom, is a moral imperative. Of course, this imperative is not imposed from the outside, but rather arises out of the very nature of the state of being.

\begin{table}[h]\centering\small
\begin{tabular}{cc}\toprule
\textbf{Kingdom} & \textbf{Imperative} \\\midrule Mineral & To be \\ Vegetable & To live \\ Animal & To be conscious \\ Human & To be rational \\ Divine & To be perfect \\\bottomrule
\end{tabular}
\label{tab:FiveStatesBeing}
\caption{States of being.}
\end{table}

That is the exoteric viewpoint. Analogous to it, these kingdoms correspond to the various strata of the soul, as we have indicated many times; that is the esoteric understanding of the schema. Various worldviews can then be categorized by combining one or more of the kingdoms into one. Some examples follow:

\begin{itemize}


\item Exoteric religion sees the Divine as distinct from creation, which comprises the other four; i.e., the dualism of the Kingdom of God and the Kingdom of this world.

\item Taoism distinguishes between the Triad of Heaven (Divine), the Human, and the World (the other three).

\item Scientific materialism only recognizes the mineral kingdom and understands life as complex crystals.

\item Naturalism regards the human state as part of the animal kingdom and is unaware of the Divine.
\end{itemize}


\paragraph{Actualization}


In Guenon's terms, we see the two tasks of salvation and liberation in this schema. At the human level, a man's task is to actualize all the possibilities available to him qua man. This involves bringing the rational into manifestation. Obviously, by ``rational'' we do not mean its modern corruptions as mere critical thinking or empty logical analysis. Rational means to be in conformity with the Logos, both exoterically as the structure of the cosmos and esoterically as the Logos is born in consciousness.

The actualization of the Divine, or \emph{theosis}, is the actualization of all possibilities. Solovyov describes it this way:

\begin{quotex}
The Kingdom of God consists of men who have ceased to be merely human and form part of a new and higher plane of existence in which their purely human ends become the means and instruments for another final purpose.

The stone exists; the plant exists and is living; the animal lives and is conscious of its life in its concrete states. Man understands the meaning of life according to ideas; the Sons of God actively realize this meaning or the perfect morel order in all things to the end.
\end{quotex}


Of course, the Sons of God are the ``elite'' as Rene Guenon defines them. The elite are not part of the kingdom of man, specifically, they are not strong men or rulers on the human plane. Rather, the Sons of God transcend the human, all-too-human, plane.

This is clearly just a cursory summary and we will develop the consequences more fully in subsequent posts. Meanwhile, I am sure that this will provide a rich source of meditation for the reader who is motivated to grasp these states of being fully in his own consciousness.


\begin{center}* * *\end{center}

\begin{footnotesize}
\begin{sffamily}
\texttt{Matt on 2014-02-21 at 14:29 said:}

``At the human level, a man's task is to actualize all the possibilities available to him qua man.''

From a Thomistic perspective, Guenon's possibilities of man can be referred to as the virtues -- the expression of the human essence as a being who can distinguish the good from evil, truth from error and keep the good and true while avoiding that which is false and evil. I think a lot of people that first discover Guenon's works make the error in thinking that the actualization of all the possibilites of man means ``engage in anything that an individual can engage in''; those who are young are probably more prone to that error (I know I was). That is why engaging in an evil is not a real possibility of man, as that would be a privation/defect, not a possibility/power available to man as man.

I apologize if that all seems redundant; but I thought it could be helpful to new readers of Guenon.

Looking forward to the subsequent posts.

\hfill

\texttt{White Rabbit on 2014-02-21 at 18:45 said:}

Am I right in taking away that actualizing one's possibilities is the same as living in accordance with one's own nature or ``being oneself''?

\hfill

\texttt{Cologero on 2014-02-22 at 10:02 said:}

Of course, you are correct, White Rabbit. However, ``living in accordance with one's own nature'' cannot be taken for granted because our own nature is opaque to us. It is something we have forgotten and our task it to remember it through self-knowledge. This first requires a process of purgation, in which we discard all our false self-images.

Then there is yet a stage beyond that, suitable for the intrepid. And that is to transcend the human state altogether.

\hfill

\texttt{George on 2014-02-23 at 16:46 said:}

How does one begin the process of purgation? Is this where disconnection from our sinfulness through the sacraments such as confession may be necessary?

\hfill

\texttt{Cologero on 2014-02-23 at 17:08 said:}

It is necessary, George. As the recent letter from Guenon that we translated here pointed out, participation in the rites is proper; ditto for Tomberg. Nevertheless, we are interested in the esoteric path. So perhaps a re-reading of The Divine Comedy would help, provided it is interpreted as a description as one's own tendencies (remember \textit{Berdyaev on the noble quality}\footnote{\url{http://www.gornahoor.net/?p=7134}}?). Then the path through hell is understood as a journey through one's own states of being.

God willing, the next blog after Gornahoor is complete will be devoted to that topic.

\hfill

\texttt{Paulo on 2014-02-26 at 12:07 said:}

``Rational means to be in conformity with the Logos, both exoterically as the structure of the cosmos and esoterically as the Logos is born in consciousness.''

To''reach''this state,it is not a matter of doing something,surrender is the case,even accepting the desire to be something special and transcending him!

\hfill

\texttt{White Rabbit on 2014-03-22 at 02:55 said:}

Cologero: does sex need to be `discovered', or is that an exception?

\hfill

\texttt{Cologero on 2014-03-23 at 10:00 said:}

White Rabbit: I think you posted this comment to the wrong blog by mistake ... I, too, have been victimized by my own ``fat fingers''. I don't see the context in this particular post that mentions neither ``sex'' nor ``discovering''. Probably it needs to be ``undiscovered'' as in \textit{The Test}\footnote{\url{http://www.gornahoor.net/?p=70}} described by Miguel Serrano.

Nevertheless, I am letting the comment slip through since we have been building up that topic starting with \textit{Interview with Aphrodite}\footnote{\url{http://www.gornahoor.net/?p=7197}}. We will follow this up with some more translations from Evola on the whole Man/Woman thing, some insights from Tomberg, and finally come to a conclusion about the ``Divine Feminine''.

Thank you for your support and interest.

\hfill

\texttt{White Rabbit on 2014-03-23 at 15:52 said:}

Well, I was thinking that gender is an aspect of one's nature that is already known.

Thank you for the answers.

\hfill

\texttt{Boreas on 2020-02-19 at 13:21 said:}

There is blessing in not being given the right to choose. I've been for years torn between conflicting viewpoints and theories and it can be very stressful and consuming.

Will there come a day when the Sons of God rule the earth or is the wish for it longing for some never to be realized utopia?

\end{sffamily}
\end{footnotesize}

